% Generated draft: requires mathematical and source review.
\documentclass[11pt]{article}
\usepackage[margin=25mm]{geometry}
\usepackage{fontspec}
\setmainfont{Latin Modern Roman}
\usepackage{amsmath,amssymb,mathtools,unicode-math}
\setmathfont{Latin Modern Math}
\usepackage{enumitem,needspace,xurl,hyperref}
\setlength{\parindent}{0pt}
\setlength{\parskip}{5pt}
\setlength{\emergencystretch}{3em}
\newcommand{\N}{\mathbb{N}}
\newcommand{\Z}{\mathbb{Z}}
\newcommand{\Q}{\mathbb{Q}}
\newcommand{\R}{\mathbb{R}}
\newcommand{\C}{\mathbb{C}}
\newcommand{\F}{\mathbb{F}}
\newcommand{\A}{\mathbb{A}}
\newcommand{\PP}{\mathbb{P}}
\newcommand{\E}{\mathbb{E}}
\newcommand{\eps}{\varepsilon}
\newcommand{\dd}{\,\mathrm{d}}
\begin{document}
\section*{UnsolvedMath ID: 701750}
\textbf{Source problem code:} catalogue entry
\subsection{Definitions and mathematical statement}
Let \(p\) be a prime and let \(\F_p\) denote the field with \(p\) elements. A pre-Lie algebra over \(\F_p\) is an \(\F_p\)-vector space \(A\) equipped with an \(\F_p\)-bilinear product satisfying
\[
(xy)z-x(yz)=(yx)z-y(xz)\qquad(x,y,z\in A).
\]
Define \(L_1(A)=A\) and
\[
L_{r+1}(A)=\operatorname{span}_{\F_p}\{xy:x\in A,\ y\in L_r(A)\}\qquad(r\geq 1).
\]
The algebra is left nilpotent if \(L_n(A)=\{0\}\) for some integer \(n\geq 1\); thus the products in this definition are nested to the right, as in the source.

A left brace is a set \(A\) with operations \(+\) and \(\circ\) such that \((A,+)\) is an abelian group, \((A,\circ)\) is a group, and
\[
a\circ(b+c)+a=a\circ b+a\circ c\qquad(a,b,c\in A).
\]
Write
\[
a*b=a\circ b-a-b.
\]
An \(\F_p\)-brace is a left brace whose additive group is equipped with an \(\F_p\)-vector-space structure satisfying
\[
a*(\alpha b)=\alpha(a*b)\qquad(a,b\in A,\ \alpha\in\F_p).
\]

\begin{enumerate}
\item[(a)] Does there exist a function \(B:\N\to\N\), with \(B(k)\geq 2\), such that for every \(k\in\N\), every prime \(p>B(k)\), and every \(\F_p\)-brace \(A\) of cardinality \(p^k\), the operation
\[
a\cdot b=-\sum_{i=0}^{p-2}(2^i)^{-1}\bigl((2^i a)*b\bigr)\qquad(a,b\in A)
\]
makes the same vector space \(A\) a pre-Lie algebra? Here \(2^i a\) denotes scalar multiplication and \((2^i)^{-1}\) is the inverse in \(\F_p\). The threshold formulation expresses the source phrase ``\(p\) sufficiently large relative to \(k\)''; no numerical bound is supplied. Taking \(B(k)\geq 2\) ensures that the displayed inverses exist.
\item[(b)] For every \(k\in\N\) and every prime \(p>2^k\), is there a bijective correspondence between \(\F_p\)-braces of cardinality \(p^k\) and left nilpotent pre-Lie algebras over \(\F_p\) of cardinality \(p^k\)? The source does not specify whether this correspondence concerns structures on a fixed underlying vector space, isomorphism classes, or a correspondence with additional compatibility properties. It also does not explicitly identify the correspondence with the operation in part~(a). These aspects of the intended target remain unresolved.
\end{enumerate}

The source includes a comment dated 2025: if
\[
2\equiv\xi^{p^{n-1}}\pmod p,
\]
where \(\xi\) is a primitive root modulo \(p\), and if \(A\) is strongly nilpotent of nilpotency index less than \(p\), then the conclusion of part~(a) follows from the cited 2022 work of A. Smoktunowicz. The parameter \(n\) and the meaning and indexing convention of strong nilpotency are not defined in the supplied record, so this comment cannot be made into a fully specified auxiliary assertion without further source clarification. The comment also states that the result was not known without strong nilpotency; this is retained as a dated source statement.

The supplied later discussion labels part~(a) an affirmative but unverified AI solution claim and labels part~(b) solved affirmatively. These are supplied source assessments, not mathematical conclusions verified in this draft. The record describes the questions as seeking an analogue of the Lazard correspondence and states that an affirmative answer to either would have consequences for set-theoretic solutions of the Yang--Baxter equation and for Hopf--Galois extensions; no precise consequences are supplied.
\subsection{Short English statement}
The questions ask whether a weighted sum of brace products always produces a pre-Lie algebra when \(p\) is sufficiently large relative to \(k\), and whether braces and left nilpotent pre-Lie algebras of size \(p^k\) correspond bijectively when \(p>2^k\). The intended meaning of correspondence and parts of the 2025 comment need clarification. The supplied record labels part~(b) solved and reports an unverified affirmative AI claim for part~(a); neither assessment is independently verified here.
\subsection{Sources}
\begin{itemize}
\item Local UnsolvedMath registry, permanent ID \(701750\), exact source code \texttt{catalogue entry}, title \emph{Kourovka Notebook Problem 20.92}. The supplied record attributes the problem to A. Smoktunowicz, proposed in 2022, and labels its overall status \texttt{partially\_solved}.
\item \emph{The Kourovka Notebook (Unsolved Problems in Group Theory)}, 20th issue (2022), problem 20.92; also the supplied 21st-issue edition (2026), arXiv:\texttt{1401.0300v48}. The registry describes the collection as edited by E. I. Khukhro and V. D. Mazurov.
\item W. Rump (2014), cited by the source for the idea of a connection between braces and pre-Lie algebras. No title, venue, or link is supplied.
\item A. Smoktunowicz, \emph{Adv. Math.} \textbf{409}, part B (2022), 108683. The source's 2025 comment attributes the qualified strongly nilpotent case of part~(a) to this work. No title or link is supplied.
\item M. Kinyon and J. Urban, preprint (2026), \url{https://kourovkanotebookorg.wordpress.com/wp-content/uploads/2026/09/kourovka-mathematics.pdf}. The supplied discussion explicitly identifies the affirmative claim for part~(a) as an unverified AI solution.
\item S. Trappeniers, preprint (2024), \url{https://arxiv.org/abs/2406.02475}. The supplied discussion labels part~(b) solved affirmatively and attributes that assessment to this reference.
\item The supplied status review is dated \texttt{2026-10-16}, identifies its evidence as \texttt{maintained\_source}, and refers to the October 2026 21st issue of the Notebook. This is supplied dated curation, not an independently verified literature assessment. The references above have not been visited or checked for this draft.
\end{itemize}
\end{document}
